\section{Curved objects and tessellation}
\label{sec:curved}
\here{1}
Curved surfaces are drawn as polygons. The silhouette error of a sphere with $s$ slices is the sagitta of one chord;
the triangle count of \code{createSphere(s, k)} is $2s(k-1)$. Halving the error costs about twice the triangles.
\begin{equation}
\varepsilon=r\Big(1-\cos\frac{\pi}{s}\Big)\approx\frac{r\pi^2}{2s^2},\qquad N_\triangle=2s(k-1)
\label{eq:tess}
\end{equation}
\begin{figure}[tbp]
\centering
\plot{p_tessellation}{\linewidth}
\caption{(a) Silhouette error against triangle count (Eq.~\eqref{eq:tess}); the dashed line is the small-angle
approximation. (b--d) The mesh \code{createSphere(s,\,s/2)} builds, shaded in the key colour; the dashed circle is the
true outline. The scene's 36 slices give 0.38\,\% of the radius with 1224 triangles.}
\label{fig:tess}
\end{figure}
Smooth shading hides the facets inside the outline but not the polygonal silhouette (Figure~\ref{fig:shading}): only
more triangles fix the edge.

\textbf{Curved objects in the scene.} Barrels, masts and heads are spheres and cylinders; tyres and muzzle rings are
tori; helmets, the radome and the bunker's earth berm are hemispheres; sandbags are capsules; the concertina wire is
made of helices and the rotor blur discs are flat cylinders. The soldiers' walking paths and the patrol routes are
curves too (centripetal Catmull--Rom splines, Section~\ref{sec:paths}). The CC0 trees carry their own detail levels:
three simplified meshes and an octahedral impostor far away.

\textbf{Limitation.} The primitives have no level of detail: all \nSphereInst{} sphere parts use \nSphereTri{}
triangles each (about \sphereTriTotal{} million triangles), and a distant sphere costs the same as a close one. Choosing
$s$ from the projected size (for example $s\propto\sqrt{\text{pixels}}$) would cut the sphere triangles by a large
factor in the overview shots.
